\documentclass[10pt,a4paper]{amsart}
\usepackage[margin=19mm]{geometry}
\usepackage{amsmath,amssymb,mathtools}
\usepackage[scr=rsfs,cal=euler]{mathalfa}
\usepackage{xfrac}
\usepackage[unicode,hidelinks,bookmarks=false]{hyperref}
\newcommand{\ie}{i.e.\ }
\newcommand{\cf}{cf.\ }
\newcommand{\resp}{respectively\ }
\newcommand{\Subsection}{\subsection}
\providecommand{\nobreakdash}{\nobreak\hbox{-}}
\setlength{\emergencystretch}{2em}
\multlinegap0pt
\usepackage{amssymb}
\newtheorem{prop}{Proposition}
\DeclareMathOperator{\C}{\mathbb{C}}
\DeclareMathOperator{\Q}{\mathbb{Q}}
\DeclareMathOperator{\Z}{\mathbb{Z}}
\DeclareMathOperator{\sF}{\mathcal{F}}
\title{Representation Rings of Fusion Systems and Brauer Characters}
\author{Thomas Lawrence}
\date{Source: arXiv:2409.03007v3 (CC BY 4.0)}
\begin{document}
\maketitle

\noindent\emph{Source and licence.} Representation Rings of Fusion Systems and Brauer Characters. Version arXiv:2409.03007v3 (CC BY 4.0), distributed by arXiv under the Creative Commons Attribution 4.0 International licence, http://creativecommons.org/licenses/by/4.0/, as recorded in the arXiv licence metadata for this exact version and shown on the abstract page https://arxiv.org/abs/2409.03007v3. Full licence notice: \texttt{licenses/arxiv-2409.03007-CC-BY-4.0.txt}.\par\smallskip

\noindent\textit{Editorial scope.} Counted unit: the article's prop character (source0.tex lines 474--478). The article's complete original proof of it is delivered with it (source0.tex lines 479--493, one \texttt{proof} environment of 15 lines). No mathematics is written, repaired or re-derived: the statement and the proof are the article's own lines, in the article's own notation, and no retained line is substituted or re-flowed.  No further source-local context is needed: every cross-reference of every retained line resolves inside the two delivered spans.  Nothing was omitted from any retained span, so the package carries no omission record.  No retained line contains a citation, so the wrapper carries no bibliography and no bibliography entry.  No retained line contains a figure environment, an image-inclusion command or drawing code, so no image is delivered and the package declares no asset dependency.  Editorial formatting only, in addition: an AMS \texttt{amsart} preamble that restates the article's own declarations and macros used by the retained lines, namely the environments \texttt{prop} on separate counters, together with the standard packages those lines need.  The retained environments print in this document as Proposition 1 in order of appearance, whereas the article prints the same items with the numbering of its own full document.  The complete native source of the article as distributed in the arXiv e-print for this exact version is bundled unchanged as \texttt{sources/164766/source0.tex}.
\begin{prop}\label{regular character}
Let $\sF = \sF_S(G)$ with $S$ not necessarily Sylow in $G$, and $\rho_S, \rho_G$ be the regular characters of $S$ and $G$. Then for any basis $B$ of $R_{\C}(\sF)$ we have
\[\rho_S = \sum_{\psi \in B} \frac{\Phi_{\psi}(1)}{[G:S]}\psi, \quad \rho_G = \sum_{\psi \in B} \psi(1)\Phi_{\psi}.\]
Furthermore, $\Phi_{\psi}(1)/[G:S]$ is an integer.
\end{prop}

\begin{proof}
We have
\begin{align*}
\mathrm{Res}^G_S(\rho_G) &= \sum_{\chi \in \text{Irr}(G)} \chi(1)\mathrm{Res}^G_S(\chi) = \sum_{\chi \in \text{Irr}(G)}\left(\chi(1)\sum_{\psi \in B} d_{\chi\phi}\psi\right) \\ 
&= \sum_{\psi \in B} \left(\psi \sum_{\chi \in \text{Irr}(G)} d_{\chi\psi}\chi(1)\right) = \sum_{\psi \in B} \psi\Phi_{\psi}(1)
\end{align*}
and because $\mathrm{Res}^G_S(\rho_G) = [G:S]\rho_S$ we have $\rho_S = \displaystyle \sum_{\psi \in B} \frac{\Phi_{\psi}(1)}{[G:S]}\psi$. Conversely, we have
{\allowdisplaybreaks\begin{align*}
\rho_G &= \sum_{\chi \in \text{Irr}(G)} \chi(1)\chi = \sum_{\chi \in \text{Irr}(G)}\left(\chi\sum_{\psi \in B} d_{\chi\phi}\psi(1)\right) \\ 
&= \sum_{\psi \in B} \left(\psi(1) \sum_{\chi \in \text{Irr}(G)} d_{\chi\psi}\chi\right) = \sum_{\psi \in B} \psi(1)\Phi_{\psi}.
\end{align*}}
Concluding the proof of the first statement, and we now show that $\frac{\Phi_{\psi}(1)}{[G:S]} \in \Z$. $\rho_S \in R_{\C}(\sF)$ implies there exist a unique set $\{a_{\psi}\}_{\psi \in B} \subset \Z$ such that $\rho_S = \sum_{\psi \in B} a_{\psi}\psi$. Because $\Q$ is a flat $\Z$-module, $B$ is a $\Q$-basis for $\Q \otimes R_{\C}(\sF)$. However, this introduces a linear relation 
\[0 = \rho_S-\rho_S = \sum_{\psi \in B}\left(a_{\psi}-\frac{\Phi_{\psi}(1)}{[G:S]}\right)\psi\]
in $\Q \otimes R_{\C}(\sF)$, which is a contradiction unless $a_{\psi} = \frac{\Phi_{\psi}(1)}{[G : S]}$. Hence $\frac{\Phi_{\psi}(1)}{[G:S]} \in \Z$.
\end{proof}

\end{document}
